% Part: first-order-logic
% Chapter: tableaux
% Section: provability-consistency

\documentclass[../../../include/open-logic-section]{subfiles}

\begin{document}

\iftag{FOL}
      {\olfileid[id]{fol}{tab}{prv}}
      {\olfileid[id]{pl}{tab}{prv}}

\olsection{\usetoken{S}{derivability} dan Konsistensi}

Sekarang kita akan membuktikan sejumlah sifat relasi !!{derivability}.
Sifat-sifat tersebut menarik secara mandiri, tetapi masing-masing akan
berperan dalam pembuktian teorema kelengkapan.

\begin{prop}\ollabel{prop:provability-contr}
  Jika $\Gamma \Proves !A$ dan $\Gamma \cup \{!A\}$ tak konsisten,
  maka $\Gamma$ tak konsisten.
\end{prop}

\begin{proof}
Terdapat himpunan berhingga $\Gamma_0 = \{!B_1, \dots, !B_n\}$ dan
$\Gamma_1 =\{!C_1, \dots, !C_m\} \subseteq \Gamma$ sedemikian sehingga
  \begin{align*}
    \{\sFmla{\False}{!A}, &
    \sFmla{\True}{!B_1}, \dots, \sFmla{\True}{!B_n}\} \\
    \{\sFmla{\True}{!A}, &
    \sFmla{\True}{!C_1}, \dots, \sFmla{\True}{!C_m}\}
  \end{align*}
  memiliki !!{tableau} tertutup. Dengan menggunakan aturan \Cut{} pada
  $!A$, kita dapat menggabungkannya menjadi satu !!{tableau} tertutup
  yang menunjukkan bahwa $\Gamma_0 \cup \Gamma_1$ tak konsisten. Karena
  $\Gamma_0 \subseteq \Gamma$ dan $\Gamma_1 \subseteq \Gamma$, maka
  $\Gamma_0 \cup \Gamma_1 \subseteq \Gamma$; jadi, $\Gamma$ tak
  konsisten.
\end{proof}

\begin{prop}
\ollabel{prop:prov-incons}
$\Gamma \Proves !A$ jika dan hanya jika $\Gamma \cup \{\lnot !A\}$
tak konsisten.
\end{prop}

\begin{proof}
Pertama, andaikan $\Gamma \Proves !A$, yakni terdapat !!{tableau}
tertutup untuk
\[
\{\sFmla{\False}{!A},
\sFmla{\True}{!B_1}, \dots, \sFmla{\True}{!B_n}\}
\]
Dengan menggunakan aturan $\TRule{\True}{\lnot}$, hasil ini dapat diubah
menjadi !!{tableau} tertutup untuk
\[
\{\sFmla{\True}{\lnot !A},
\sFmla{\True}{!B_1}, \dots, \sFmla{\True}{!B_n}\}.
\]

Sebaliknya, jika terdapat !!{tableau} tertutup untuk himpunan kedua,
kita dapat mengubahnya menjadi !!{tableau} tertutup untuk himpunan
pertama dengan menghapus setiap formula yang dihasilkan dari penerapan
\TRule{\True}{\lnot} pada asumsi pertama
$\sFmla{\True}{\lnot !A}$ beserta asumsi tersebut, lalu menambahkan
asumsi $\sFmla{\False}{!A}$. Jika sebelumnya suatu cabang tertutup karena
memuat kesimpulan penerapan \TRule{\True}{\lnot} pada
$\sFmla{\True}{\lnot !A}$, yakni $\sFmla{\False}{!A}$, cabang yang
bersesuaian dalam !!{tableau} baru juga tertutup. Jika suatu cabang dalam
tableau lama tertutup karena memuat asumsi
$\sFmla{\True}{\lnot !A}$ serta $\sFmla{\False}{\lnot !A}$, kita dapat
mengubahnya menjadi cabang tertutup dengan menerapkan
$\TRule{\False}{\lnot}$ pada $\sFmla{\False}{\lnot !A}$ untuk memperoleh
$\sFmla{\True}{!A}$. Cabang tersebut tertutup karena kita telah
menambahkan $\sFmla{\False}{!A}$ sebagai asumsi.
\end{proof}

\begin{prob}
Buktikan bahwa $\Gamma \Proves \lnot !A$ jika dan hanya jika
$\Gamma \cup \{!A\}$ tak konsisten.
\end{prob}

\begin{prop}\ollabel{prop:explicit-inc}
  Jika $\Gamma \Proves !A$ dan $\lnot !A \in \Gamma$, maka $\Gamma$
  tak konsisten.
\end{prop}

\begin{proof}
  Andaikan $\Gamma \Proves !A$ dan $\lnot !A \in \Gamma$. Maka terdapat
  $!B_1$, \dots, $!B_n \in \Gamma$ sedemikian sehingga \ 
  \[
  \{\sFmla{\False}{!A}, \sFmla{\True}{!B_1}, \dots,
  \sFmla{\True}{!B_n}\}
  \]
  memiliki tableau tertutup. Ganti asumsi \sFmla{\False}{!A} dengan
  \sFmla{\True}{\lnot !A}, lalu sisipkan kesimpulan penerapan
  \TRule{\True}{\lnot} pada \sFmla{\True}{\lnot !A} setelah semua
  asumsi. Setiap !!{sentence} dalam !!{tableau} yang sebelumnya
  dijustifikasi dengan mengacu pada baris~$1$ dalam !!{tableau} lama kini
  dijustifikasi dengan mengacu pada baris~$n+2$. Jadi, jika !!{tableau}
  lama tertutup, tableau baru pun tertutup. Hal ini menunjukkan bahwa
  $\Gamma$ tak konsisten, karena semua asumsinya
  berada dalam~$\Gamma$.
\end{proof}

\begin{prop}\ollabel{prop:provability-exhaustive}
  Jika $\Gamma \cup \{!A\}$ dan $\Gamma \cup \{\lnot !A\}$ keduanya
  tak konsisten, maka $\Gamma$ tak konsisten.
\end{prop}

\begin{proof}
  Jika terdapat $!B_1$, \dots, $!B_n \in \Gamma$ dan $!C_1$, \dots,
  $!C_m \in \Gamma$ sedemikian sehingga
  \begin{align*}
    \{\sFmla{\True}{!A}, &
    \sFmla{\True}{!B_1}, \dots, \sFmla{\True}{!B_n}\} \text{ dan}\\
    \{\sFmla{\True}{\lnot !A}, &
    \sFmla{\True}{!C_1}, \dots, \sFmla{\True}{!C_m}\}
  \end{align*}
  keduanya memiliki !!{tableau} tertutup, kita dapat membangun satu
  !!{tableau} gabungan yang menunjukkan bahwa $\Gamma$ tak konsisten.
  Gunakan $\sFmla{\True}{!B_1}$, \dots, $\sFmla{\True}{!B_n}$ bersama
  $\sFmla{\True}{!C_1}$, \dots, $\sFmla{\True}{!C_m}$ sebagai asumsi,
  lalu terapkan aturan \Cut{}. Penerapan ini menghasilkan dua cabang:
  yang satu dimulai dengan $\sFmla{\True}{!A}$, sedangkan yang lain
  dengan $\sFmla{\False}{!A}$.

  Pada sisi kiri, tambahkan bagian !!{tableau} pertama yang berada di
  bawah asumsi-asumsinya. Di sini, setiap penerapan aturan tetap benar
  karena semua asumsi !!{tableau} pertama, termasuk
  $\sFmla{\True}{!A}$, tersedia. Jadi, setiap cabang di bawah
  $\sFmla{\True}{!A}$ tertutup.

  Pada sisi kanan, tambahkan bagian !!{tableau} kedua yang berada di
  bawah asumsi-asumsinya, dengan menghapus hasil setiap penerapan
  $\TRule{\True}{\lnot}$ pada $\sFmla{\True}{\lnot !A}$. Kesimpulan
  penerapan $\TRule{\True}{\lnot}$ pada $\sFmla{\True}{\lnot !A}$ ialah
  $\sFmla{\False}{!A}$, yang tetap tersedia sebagai kesimpulan aturan
  \Cut{} pada sisi kanan !!{tableau} gabungan.

  Jika suatu cabang dalam tableau kedua tertutup karena memuat asumsi
  $\sFmla{\True}{\lnot !A}$ (yang tidak lagi muncul sebagai asumsi dalam
  !!{tableau} gabungan) serta $\sFmla{\False}{\lnot !A}$, kita dapat
  menerapkan $\TRule{\False}{\lnot}$ pada
  $\sFmla{\False}{\lnot !A}$ untuk memperoleh $\sFmla{\True}{!A}$.
  Sekarang cabang yang bersesuaian dalam !!{tableau} gabungan juga
  tertutup karena memuat kesimpulan sisi kanan aturan \Cut{}, yakni
  $\sFmla{\False}{!A}$. Jika suatu cabang dalam !!{tableau} kedua
  tertutup karena alasan lain, cabang yang bersesuaian dalam !!{tableau}
  gabungan juga tertutup, sebab setiap !!{signed formula} selain
  $\sFmla{\True}{\lnot !A}$ yang muncul pada cabang dalam !!{tableau}
  kedua yang lama juga muncul pada cabang yang bersesuaian dalam
  !!{tableau} gabungan.
\end{proof}

\end{document}
